% =====================================================================
%  Capitol / Secțiune: Testarea automată a aplicației
%  Se poate include în main.tex prin: % =====================================================================
%  Capitol / Secțiune: Testarea automată a aplicației
%  Se poate include în main.tex prin: % =====================================================================
%  Capitol / Secțiune: Testarea automată a aplicației
%  Se poate include în main.tex prin: % =====================================================================
%  Capitol / Secțiune: Testarea automată a aplicației
%  Se poate include în main.tex prin: \input{testare}
% =====================================================================

\section{Testarea automată a sistemului}
\label{sec:testare}

\subsection{Rolul și obiectivele testării}

Testarea automată reprezintă o componentă esențială a procesului de dezvoltare
a aplicației, având rolul de a verifica în mod sistematic și reproductibil
corectitudinea funcțională a fiecărei unități de cod. Într-un sistem de
prelucrare a datelor și de analiză comportamentală, în care rezultatele finale
(clasificarea participanților, estimarea gradului de dizabilitate, metricile de
evaluare) depind de un lanț lung de transformări numerice, o singură eroare
introdusă într-o funcție de bază se poate propaga silențios până la concluziile
științifice. Suita de teste are tocmai rolul de a interpune o barieră de
verificare automată între implementare și utilizare.

Obiectivele urmărite prin construirea suitei de teste au fost următoarele:

\begin{itemize}
  \item \textbf{Verificarea corectitudinii funcționale} a fiecărei funcții
        publice din modulele aplicației, pe baza unor date de intrare cu
        rezultat cunoscut a priori.
  \item \textbf{Protejarea împotriva regresiilor}: orice modificare ulterioară a
        codului poate fi validată prin re-rularea suitei, asigurând că
        funcționalitățile deja corecte nu sunt deteriorate.
  \item \textbf{Documentarea executabilă a comportamentului așteptat}: testele
        descriu, prin exemple concrete, contractul fiecărei funcții (intrări
        valide, valori de frontieră, cazuri degenerate), constituind o formă de
        specificație vie a codului.
  \item \textbf{Validarea robusteții} în fața cazurilor limită (date lipsă,
        traiectorii cu un singur cadru, distribuții identice, clase unice),
        care apar frecvent în datele experimentale reale.
  \item \textbf{Izolarea defectelor}: în cazul unei erori, granularitatea fină a
        testelor permite localizarea rapidă a modulului și a funcției
        responsabile.
\end{itemize}

\subsection{Metodologia de testare}

Strategia adoptată este aceea a \emph{testării unitare} (unit testing),
completată de elemente de \emph{testare comportamentală} și de verificare a
cazurilor de frontieră. Fiecare test urmează tiparul \emph{Arrange--Act--Assert}
(pregătirea datelor de intrare, apelul funcției testate, verificarea
rezultatului prin aserțiuni), ceea ce asigură lizibilitate și o corespondență
clară între intenția testului și implementarea sa.

Categoriile de teste utilizate sunt:

\begin{itemize}
  \item \textbf{Teste de corectitudine numerică} --- verifică egalitatea
        (eventual aproximativă, cu toleranță controlată) între valoarea
        calculată și cea așteptată teoretic; de exemplu, faptul că distanța
        totală parcursă este nulă pentru o traiectorie staționară.
  \item \textbf{Teste de proprietate / invarianți} --- verifică faptul că
        rezultatul respectă constrângeri structurale independente de datele
        concrete: eficiența traiectoriei este întotdeauna în intervalul
        $[0,1]$, procentele aferente grupurilor de clustering însumează
        $100\%$, vectorul de trăsături nu conține valori \texttt{NaN} sau
        infinite.
  \item \textbf{Teste de cazuri limită} (edge cases) --- traiectorii cu un
        singur cadru sau cu număr minim de cadre, distribuții complet
        separate sau identice în testele statistice, seturi de date vide.
  \item \textbf{Teste de regresie} --- introduse explicit în urma identificării
        și corectării unor defecte, pentru a garanta că acestea nu reapar.
\end{itemize}

\subsection{Implementarea suitei de teste}

\paragraph{Cadrul de testare.} Suita este implementată folosind
\texttt{pytest}, cadrul standard de facto pentru testarea în Python, ales pentru
sintaxa concisă a aserțiunilor (utilizarea directă a instrucțiunii
\texttt{assert}), pentru mecanismul de \emph{fixtures} și pentru capacitățile
avansate de colectare și raportare. Pentru comparațiile numerice se utilizează
\texttt{pytest.approx} și funcțiile \texttt{numpy.testing} (de exemplu
\texttt{assert\_array\_almost\_equal}), care permit specificarea unei toleranțe
controlate, necesară din cauza erorilor de rotunjire inerente aritmeticii în
virgulă mobilă.

\paragraph{Organizarea.} Testele sunt grupate în directorul dedicat
\texttt{tests/}, câte un fișier pentru fiecare modul al aplicației, iar în
interiorul fiecărui fișier testele sunt structurate pe clase care corespund
funcțiilor publice testate (de exemplu \texttt{TestComputeMannWhitney},
\texttt{TestPerformAgglomerativeClustering}). Această corespondență directă
între structura testelor și structura codului facilitează navigarea și
mentenanța.

\paragraph{Configurarea mediului.} Fișierul \texttt{conftest.py} configurează
calea de import (\texttt{sys.path}) astfel încât modulele din \texttt{src/} să
poată fi importate de teste, asigurând independența suitei față de modul de
lansare. Generarea graficelor cu \texttt{matplotlib} este comutată pe backend-ul
neinteractiv \texttt{Agg}, pentru ca testele să poată rula într-un mediu fără
interfață grafică (de exemplu pe un server de integrare continuă).

\paragraph{Izolarea dependențelor.} Pentru a testa fiecare unitate independent
de dependențele sale externe sau costisitoare, s-a folosit tehnica
\emph{mocking}-ului (prin \texttt{unittest.mock}). Astfel, funcțiile de
încărcare a datelor de pe disc, modulul de configurare și anumite biblioteci de
vizualizare sunt înlocuite cu obiecte simulate (\emph{mock objects}), care
returnează valori controlate. Aceasta permite testarea logicii unei funcții
fără a depinde de existența fișierelor reale de date VR, de starea sistemului de
fișiere sau de randarea efectivă a graficelor. Operațiunile de scriere a
fișierelor sunt redirecționate către directoare temporare furnizate de
fixture-ul \texttt{tmp\_path}, astfel încât testele să nu producă efecte
secundare persistente.

\paragraph{Acoperirea.} Suita cuprinde \textbf{353 de teste} care acoperă toate
modulele funcționale ale aplicației. Tabelul~\ref{tab:teste} prezintă
distribuția testelor pe module, împreună cu responsabilitatea fiecărui modul
testat.

\begin{table}[ht]
\centering
\caption{Distribuția testelor pe module}
\label{tab:teste}
\small
\begin{tabular}{l c p{6.5cm}}
\hline
\textbf{Modul testat} & \textbf{Nr. teste} & \textbf{Responsabilitate} \\
\hline
\texttt{optimize\_config}     & 55 & Optimizarea pragurilor și a ponderilor de decizie, încărcarea scorurilor, actualizarea configurației \\
\texttt{file\_management}     & 49 & Gruparea fișierelor pe participant, agregarea trăsăturilor, gestionarea fișierelor de ieșire \\
\texttt{disability\_engine}   & 40 & Determinarea numărului optim de clustere, estimarea statusului de dizabilitate, praguri derivate din date, analiza per participant \\
\texttt{feature\_extractor}   & 26 & Extragerea vectorului de trăsături comportamentale din traiectorie \\
\texttt{scenario\_comparison} & 22 & Compararea și vizualizarea scenariilor pe participanți \\
\texttt{loocv}                & 21 & Validarea încrucișată leave-one-out, praguri Youden și F1 \\
\texttt{ground\_truth\_handler} & 21 & Încărcarea și sincronizarea etichetelor de referință \\
\texttt{reports}              & 20 & Generarea rapoartelor vizuale și a dendrogramelor \\
\texttt{data\_loader}         & 20 & Parsarea vectorilor și încărcarea înregistrărilor VR \\
\texttt{threshold\_calibrator}& 18 & Calibrarea pragului de decizie pe baza scorurilor \\
\texttt{metrics}              & 18 & Calculul și exportul metricilor de evaluare \\
\texttt{mann\_whitney}        & 18 & Testul statistic Mann--Whitney U și mărimea efectului \\
\texttt{plotter\_3d}          & 15 & Vizualizarea tridimensională a clusterelor \\
\texttt{clustering}           & 10 & Clusteringul aglomerativ ierarhic \\
\hline
\textbf{Total}                & \textbf{353} & \\
\hline
\end{tabular}
\end{table}

\subsection{Rezultate}

În urma execuției complete a suitei, toate cele \textbf{353 de teste sunt
trecute cu succes} (\emph{353 passed}), fără eșecuri sau erori. Timpul total de
execuție este de aproximativ 40 de secunde, ceea ce face suita suficient de
rapidă pentru a fi rulată frecvent în timpul dezvoltării. Avertismentele rămase
(18 \emph{warnings}) provin din biblioteci externe (\texttt{numpy},
\texttt{scikit-learn}) și corespund unor situații numerice degenerate
\emph{intenționat} provocate de teste --- de exemplu calculul ariei de sub
curba ROC (AUC) atunci când există o singură clasă în datele de referință ---,
fără a indica defecte ale aplicației.

\paragraph{Defecte identificate și corectate.} Procesul de testare și-a
demonstrat utilitatea prin identificarea unui set de defecte, dintre care cele
mai relevante au fost:

\begin{itemize}
  \item În modulul de testare statistică Mann--Whitney, câmpurile booleene de
        decizie (\texttt{reject\_h0\_alpha05}) erau returnate ca tip boolean
        \texttt{numpy} în loc de tipul boolean nativ Python, ceea ce afecta
        comparațiile stricte de identitate; defectul a fost remediat prin
        conversie explicită de tip.
  \item În testele de extragere a trăsăturilor, indicii care identificau
        trăsăturile \emph{eficiența traiectoriei} și \emph{rata de pauză} nu mai
        corespundeau structurii actualizate a vectorului de trăsături (care a
        fost extins cu descriptori suplimentari precum tremorul, jerk-ul și
        asimetria direcțională); indicii au fost actualizați la pozițiile
        corecte (37, respectiv 43).
  \item Mai multe defecte erau cauzate de \emph{interferența} dintre module de
        test care modificau global tabela de module a interpretorului
        (\texttt{sys.modules}), substituind biblioteci reale cu obiecte
        simulate; izolarea acestor substituiri a eliminat dependențele de ordine
        dintre teste și a restabilit determinismul suitei.
\end{itemize}

Aceste corecții ilustrează rolul dublu al testării: pe de o parte validarea
codului de producție, pe de altă parte menținerea calității și a izolării
propriei infrastructuri de testare.

\subsection{Reproductibilitate și limitări}

Suita de teste este complet deterministă: generatoarele de numere aleatoare sunt
inițializate cu semințe fixe (\texttt{seed}), astfel încât rezultatele sunt
identice la fiecare rulare și independente de mediul de execuție. Lansarea se
face printr-o singură comandă (\texttt{pytest tests/}), iar dependențele sunt
cele declarate în mediul virtual al proiectului.

Ca limitări, suita se concentrează pe testarea unitară și de proprietate, nu
include teste de integrare \emph{end-to-end} care să parcurgă întregul flux de
prelucrare pe date reale, și nu cuantifică explicit acoperirea de cod
(\emph{code coverage}) printr-un instrument dedicat. Aceste direcții constituie
posibile extinderi viitoare, alături de adăugarea unor teste de performanță și a
integrării într-un sistem de integrare continuă (CI) care să ruleze automat
suita la fiecare modificare a codului.

% =====================================================================

\section{Testarea automată a sistemului}
\label{sec:testare}

\subsection{Rolul și obiectivele testării}

Testarea automată reprezintă o componentă esențială a procesului de dezvoltare
a aplicației, având rolul de a verifica în mod sistematic și reproductibil
corectitudinea funcțională a fiecărei unități de cod. Într-un sistem de
prelucrare a datelor și de analiză comportamentală, în care rezultatele finale
(clasificarea participanților, estimarea gradului de dizabilitate, metricile de
evaluare) depind de un lanț lung de transformări numerice, o singură eroare
introdusă într-o funcție de bază se poate propaga silențios până la concluziile
științifice. Suita de teste are tocmai rolul de a interpune o barieră de
verificare automată între implementare și utilizare.

Obiectivele urmărite prin construirea suitei de teste au fost următoarele:

\begin{itemize}
  \item \textbf{Verificarea corectitudinii funcționale} a fiecărei funcții
        publice din modulele aplicației, pe baza unor date de intrare cu
        rezultat cunoscut a priori.
  \item \textbf{Protejarea împotriva regresiilor}: orice modificare ulterioară a
        codului poate fi validată prin re-rularea suitei, asigurând că
        funcționalitățile deja corecte nu sunt deteriorate.
  \item \textbf{Documentarea executabilă a comportamentului așteptat}: testele
        descriu, prin exemple concrete, contractul fiecărei funcții (intrări
        valide, valori de frontieră, cazuri degenerate), constituind o formă de
        specificație vie a codului.
  \item \textbf{Validarea robusteții} în fața cazurilor limită (date lipsă,
        traiectorii cu un singur cadru, distribuții identice, clase unice),
        care apar frecvent în datele experimentale reale.
  \item \textbf{Izolarea defectelor}: în cazul unei erori, granularitatea fină a
        testelor permite localizarea rapidă a modulului și a funcției
        responsabile.
\end{itemize}

\subsection{Metodologia de testare}

Strategia adoptată este aceea a \emph{testării unitare} (unit testing),
completată de elemente de \emph{testare comportamentală} și de verificare a
cazurilor de frontieră. Fiecare test urmează tiparul \emph{Arrange--Act--Assert}
(pregătirea datelor de intrare, apelul funcției testate, verificarea
rezultatului prin aserțiuni), ceea ce asigură lizibilitate și o corespondență
clară între intenția testului și implementarea sa.

Categoriile de teste utilizate sunt:

\begin{itemize}
  \item \textbf{Teste de corectitudine numerică} --- verifică egalitatea
        (eventual aproximativă, cu toleranță controlată) între valoarea
        calculată și cea așteptată teoretic; de exemplu, faptul că distanța
        totală parcursă este nulă pentru o traiectorie staționară.
  \item \textbf{Teste de proprietate / invarianți} --- verifică faptul că
        rezultatul respectă constrângeri structurale independente de datele
        concrete: eficiența traiectoriei este întotdeauna în intervalul
        $[0,1]$, procentele aferente grupurilor de clustering însumează
        $100\%$, vectorul de trăsături nu conține valori \texttt{NaN} sau
        infinite.
  \item \textbf{Teste de cazuri limită} (edge cases) --- traiectorii cu un
        singur cadru sau cu număr minim de cadre, distribuții complet
        separate sau identice în testele statistice, seturi de date vide.
  \item \textbf{Teste de regresie} --- introduse explicit în urma identificării
        și corectării unor defecte, pentru a garanta că acestea nu reapar.
\end{itemize}

\subsection{Implementarea suitei de teste}

\paragraph{Cadrul de testare.} Suita este implementată folosind
\texttt{pytest}, cadrul standard de facto pentru testarea în Python, ales pentru
sintaxa concisă a aserțiunilor (utilizarea directă a instrucțiunii
\texttt{assert}), pentru mecanismul de \emph{fixtures} și pentru capacitățile
avansate de colectare și raportare. Pentru comparațiile numerice se utilizează
\texttt{pytest.approx} și funcțiile \texttt{numpy.testing} (de exemplu
\texttt{assert\_array\_almost\_equal}), care permit specificarea unei toleranțe
controlate, necesară din cauza erorilor de rotunjire inerente aritmeticii în
virgulă mobilă.

\paragraph{Organizarea.} Testele sunt grupate în directorul dedicat
\texttt{tests/}, câte un fișier pentru fiecare modul al aplicației, iar în
interiorul fiecărui fișier testele sunt structurate pe clase care corespund
funcțiilor publice testate (de exemplu \texttt{TestComputeMannWhitney},
\texttt{TestPerformAgglomerativeClustering}). Această corespondență directă
între structura testelor și structura codului facilitează navigarea și
mentenanța.

\paragraph{Configurarea mediului.} Fișierul \texttt{conftest.py} configurează
calea de import (\texttt{sys.path}) astfel încât modulele din \texttt{src/} să
poată fi importate de teste, asigurând independența suitei față de modul de
lansare. Generarea graficelor cu \texttt{matplotlib} este comutată pe backend-ul
neinteractiv \texttt{Agg}, pentru ca testele să poată rula într-un mediu fără
interfață grafică (de exemplu pe un server de integrare continuă).

\paragraph{Izolarea dependențelor.} Pentru a testa fiecare unitate independent
de dependențele sale externe sau costisitoare, s-a folosit tehnica
\emph{mocking}-ului (prin \texttt{unittest.mock}). Astfel, funcțiile de
încărcare a datelor de pe disc, modulul de configurare și anumite biblioteci de
vizualizare sunt înlocuite cu obiecte simulate (\emph{mock objects}), care
returnează valori controlate. Aceasta permite testarea logicii unei funcții
fără a depinde de existența fișierelor reale de date VR, de starea sistemului de
fișiere sau de randarea efectivă a graficelor. Operațiunile de scriere a
fișierelor sunt redirecționate către directoare temporare furnizate de
fixture-ul \texttt{tmp\_path}, astfel încât testele să nu producă efecte
secundare persistente.

\paragraph{Acoperirea.} Suita cuprinde \textbf{353 de teste} care acoperă toate
modulele funcționale ale aplicației. Tabelul~\ref{tab:teste} prezintă
distribuția testelor pe module, împreună cu responsabilitatea fiecărui modul
testat.

\begin{table}[ht]
\centering
\caption{Distribuția testelor pe module}
\label{tab:teste}
\small
\begin{tabular}{l c p{6.5cm}}
\hline
\textbf{Modul testat} & \textbf{Nr. teste} & \textbf{Responsabilitate} \\
\hline
\texttt{optimize\_config}     & 55 & Optimizarea pragurilor și a ponderilor de decizie, încărcarea scorurilor, actualizarea configurației \\
\texttt{file\_management}     & 49 & Gruparea fișierelor pe participant, agregarea trăsăturilor, gestionarea fișierelor de ieșire \\
\texttt{disability\_engine}   & 40 & Determinarea numărului optim de clustere, estimarea statusului de dizabilitate, praguri derivate din date, analiza per participant \\
\texttt{feature\_extractor}   & 26 & Extragerea vectorului de trăsături comportamentale din traiectorie \\
\texttt{scenario\_comparison} & 22 & Compararea și vizualizarea scenariilor pe participanți \\
\texttt{loocv}                & 21 & Validarea încrucișată leave-one-out, praguri Youden și F1 \\
\texttt{ground\_truth\_handler} & 21 & Încărcarea și sincronizarea etichetelor de referință \\
\texttt{reports}              & 20 & Generarea rapoartelor vizuale și a dendrogramelor \\
\texttt{data\_loader}         & 20 & Parsarea vectorilor și încărcarea înregistrărilor VR \\
\texttt{threshold\_calibrator}& 18 & Calibrarea pragului de decizie pe baza scorurilor \\
\texttt{metrics}              & 18 & Calculul și exportul metricilor de evaluare \\
\texttt{mann\_whitney}        & 18 & Testul statistic Mann--Whitney U și mărimea efectului \\
\texttt{plotter\_3d}          & 15 & Vizualizarea tridimensională a clusterelor \\
\texttt{clustering}           & 10 & Clusteringul aglomerativ ierarhic \\
\hline
\textbf{Total}                & \textbf{353} & \\
\hline
\end{tabular}
\end{table}

\subsection{Rezultate}

În urma execuției complete a suitei, toate cele \textbf{353 de teste sunt
trecute cu succes} (\emph{353 passed}), fără eșecuri sau erori. Timpul total de
execuție este de aproximativ 40 de secunde, ceea ce face suita suficient de
rapidă pentru a fi rulată frecvent în timpul dezvoltării. Avertismentele rămase
(18 \emph{warnings}) provin din biblioteci externe (\texttt{numpy},
\texttt{scikit-learn}) și corespund unor situații numerice degenerate
\emph{intenționat} provocate de teste --- de exemplu calculul ariei de sub
curba ROC (AUC) atunci când există o singură clasă în datele de referință ---,
fără a indica defecte ale aplicației.

\paragraph{Defecte identificate și corectate.} Procesul de testare și-a
demonstrat utilitatea prin identificarea unui set de defecte, dintre care cele
mai relevante au fost:

\begin{itemize}
  \item În modulul de testare statistică Mann--Whitney, câmpurile booleene de
        decizie (\texttt{reject\_h0\_alpha05}) erau returnate ca tip boolean
        \texttt{numpy} în loc de tipul boolean nativ Python, ceea ce afecta
        comparațiile stricte de identitate; defectul a fost remediat prin
        conversie explicită de tip.
  \item În testele de extragere a trăsăturilor, indicii care identificau
        trăsăturile \emph{eficiența traiectoriei} și \emph{rata de pauză} nu mai
        corespundeau structurii actualizate a vectorului de trăsături (care a
        fost extins cu descriptori suplimentari precum tremorul, jerk-ul și
        asimetria direcțională); indicii au fost actualizați la pozițiile
        corecte (37, respectiv 43).
  \item Mai multe defecte erau cauzate de \emph{interferența} dintre module de
        test care modificau global tabela de module a interpretorului
        (\texttt{sys.modules}), substituind biblioteci reale cu obiecte
        simulate; izolarea acestor substituiri a eliminat dependențele de ordine
        dintre teste și a restabilit determinismul suitei.
\end{itemize}

Aceste corecții ilustrează rolul dublu al testării: pe de o parte validarea
codului de producție, pe de altă parte menținerea calității și a izolării
propriei infrastructuri de testare.

\subsection{Reproductibilitate și limitări}

Suita de teste este complet deterministă: generatoarele de numere aleatoare sunt
inițializate cu semințe fixe (\texttt{seed}), astfel încât rezultatele sunt
identice la fiecare rulare și independente de mediul de execuție. Lansarea se
face printr-o singură comandă (\texttt{pytest tests/}), iar dependențele sunt
cele declarate în mediul virtual al proiectului.

Ca limitări, suita se concentrează pe testarea unitară și de proprietate, nu
include teste de integrare \emph{end-to-end} care să parcurgă întregul flux de
prelucrare pe date reale, și nu cuantifică explicit acoperirea de cod
(\emph{code coverage}) printr-un instrument dedicat. Aceste direcții constituie
posibile extinderi viitoare, alături de adăugarea unor teste de performanță și a
integrării într-un sistem de integrare continuă (CI) care să ruleze automat
suita la fiecare modificare a codului.

% =====================================================================

\section{Testarea automată a sistemului}
\label{sec:testare}

\subsection{Rolul și obiectivele testării}

Testarea automată reprezintă o componentă esențială a procesului de dezvoltare
a aplicației, având rolul de a verifica în mod sistematic și reproductibil
corectitudinea funcțională a fiecărei unități de cod. Într-un sistem de
prelucrare a datelor și de analiză comportamentală, în care rezultatele finale
(clasificarea participanților, estimarea gradului de dizabilitate, metricile de
evaluare) depind de un lanț lung de transformări numerice, o singură eroare
introdusă într-o funcție de bază se poate propaga silențios până la concluziile
științifice. Suita de teste are tocmai rolul de a interpune o barieră de
verificare automată între implementare și utilizare.

Obiectivele urmărite prin construirea suitei de teste au fost următoarele:

\begin{itemize}
  \item \textbf{Verificarea corectitudinii funcționale} a fiecărei funcții
        publice din modulele aplicației, pe baza unor date de intrare cu
        rezultat cunoscut a priori.
  \item \textbf{Protejarea împotriva regresiilor}: orice modificare ulterioară a
        codului poate fi validată prin re-rularea suitei, asigurând că
        funcționalitățile deja corecte nu sunt deteriorate.
  \item \textbf{Documentarea executabilă a comportamentului așteptat}: testele
        descriu, prin exemple concrete, contractul fiecărei funcții (intrări
        valide, valori de frontieră, cazuri degenerate), constituind o formă de
        specificație vie a codului.
  \item \textbf{Validarea robusteții} în fața cazurilor limită (date lipsă,
        traiectorii cu un singur cadru, distribuții identice, clase unice),
        care apar frecvent în datele experimentale reale.
  \item \textbf{Izolarea defectelor}: în cazul unei erori, granularitatea fină a
        testelor permite localizarea rapidă a modulului și a funcției
        responsabile.
\end{itemize}

\subsection{Metodologia de testare}

Strategia adoptată este aceea a \emph{testării unitare} (unit testing),
completată de elemente de \emph{testare comportamentală} și de verificare a
cazurilor de frontieră. Fiecare test urmează tiparul \emph{Arrange--Act--Assert}
(pregătirea datelor de intrare, apelul funcției testate, verificarea
rezultatului prin aserțiuni), ceea ce asigură lizibilitate și o corespondență
clară între intenția testului și implementarea sa.

Categoriile de teste utilizate sunt:

\begin{itemize}
  \item \textbf{Teste de corectitudine numerică} --- verifică egalitatea
        (eventual aproximativă, cu toleranță controlată) între valoarea
        calculată și cea așteptată teoretic; de exemplu, faptul că distanța
        totală parcursă este nulă pentru o traiectorie staționară.
  \item \textbf{Teste de proprietate / invarianți} --- verifică faptul că
        rezultatul respectă constrângeri structurale independente de datele
        concrete: eficiența traiectoriei este întotdeauna în intervalul
        $[0,1]$, procentele aferente grupurilor de clustering însumează
        $100\%$, vectorul de trăsături nu conține valori \texttt{NaN} sau
        infinite.
  \item \textbf{Teste de cazuri limită} (edge cases) --- traiectorii cu un
        singur cadru sau cu număr minim de cadre, distribuții complet
        separate sau identice în testele statistice, seturi de date vide.
  \item \textbf{Teste de regresie} --- introduse explicit în urma identificării
        și corectării unor defecte, pentru a garanta că acestea nu reapar.
\end{itemize}

\subsection{Implementarea suitei de teste}

\paragraph{Cadrul de testare.} Suita este implementată folosind
\texttt{pytest}, cadrul standard de facto pentru testarea în Python, ales pentru
sintaxa concisă a aserțiunilor (utilizarea directă a instrucțiunii
\texttt{assert}), pentru mecanismul de \emph{fixtures} și pentru capacitățile
avansate de colectare și raportare. Pentru comparațiile numerice se utilizează
\texttt{pytest.approx} și funcțiile \texttt{numpy.testing} (de exemplu
\texttt{assert\_array\_almost\_equal}), care permit specificarea unei toleranțe
controlate, necesară din cauza erorilor de rotunjire inerente aritmeticii în
virgulă mobilă.

\paragraph{Organizarea.} Testele sunt grupate în directorul dedicat
\texttt{tests/}, câte un fișier pentru fiecare modul al aplicației, iar în
interiorul fiecărui fișier testele sunt structurate pe clase care corespund
funcțiilor publice testate (de exemplu \texttt{TestComputeMannWhitney},
\texttt{TestPerformAgglomerativeClustering}). Această corespondență directă
între structura testelor și structura codului facilitează navigarea și
mentenanța.

\paragraph{Configurarea mediului.} Fișierul \texttt{conftest.py} configurează
calea de import (\texttt{sys.path}) astfel încât modulele din \texttt{src/} să
poată fi importate de teste, asigurând independența suitei față de modul de
lansare. Generarea graficelor cu \texttt{matplotlib} este comutată pe backend-ul
neinteractiv \texttt{Agg}, pentru ca testele să poată rula într-un mediu fără
interfață grafică (de exemplu pe un server de integrare continuă).

\paragraph{Izolarea dependențelor.} Pentru a testa fiecare unitate independent
de dependențele sale externe sau costisitoare, s-a folosit tehnica
\emph{mocking}-ului (prin \texttt{unittest.mock}). Astfel, funcțiile de
încărcare a datelor de pe disc, modulul de configurare și anumite biblioteci de
vizualizare sunt înlocuite cu obiecte simulate (\emph{mock objects}), care
returnează valori controlate. Aceasta permite testarea logicii unei funcții
fără a depinde de existența fișierelor reale de date VR, de starea sistemului de
fișiere sau de randarea efectivă a graficelor. Operațiunile de scriere a
fișierelor sunt redirecționate către directoare temporare furnizate de
fixture-ul \texttt{tmp\_path}, astfel încât testele să nu producă efecte
secundare persistente.

\paragraph{Acoperirea.} Suita cuprinde \textbf{353 de teste} care acoperă toate
modulele funcționale ale aplicației. Tabelul~\ref{tab:teste} prezintă
distribuția testelor pe module, împreună cu responsabilitatea fiecărui modul
testat.

\begin{table}[ht]
\centering
\caption{Distribuția testelor pe module}
\label{tab:teste}
\small
\begin{tabular}{l c p{6.5cm}}
\hline
\textbf{Modul testat} & \textbf{Nr. teste} & \textbf{Responsabilitate} \\
\hline
\texttt{optimize\_config}     & 55 & Optimizarea pragurilor și a ponderilor de decizie, încărcarea scorurilor, actualizarea configurației \\
\texttt{file\_management}     & 49 & Gruparea fișierelor pe participant, agregarea trăsăturilor, gestionarea fișierelor de ieșire \\
\texttt{disability\_engine}   & 40 & Determinarea numărului optim de clustere, estimarea statusului de dizabilitate, praguri derivate din date, analiza per participant \\
\texttt{feature\_extractor}   & 26 & Extragerea vectorului de trăsături comportamentale din traiectorie \\
\texttt{scenario\_comparison} & 22 & Compararea și vizualizarea scenariilor pe participanți \\
\texttt{loocv}                & 21 & Validarea încrucișată leave-one-out, praguri Youden și F1 \\
\texttt{ground\_truth\_handler} & 21 & Încărcarea și sincronizarea etichetelor de referință \\
\texttt{reports}              & 20 & Generarea rapoartelor vizuale și a dendrogramelor \\
\texttt{data\_loader}         & 20 & Parsarea vectorilor și încărcarea înregistrărilor VR \\
\texttt{threshold\_calibrator}& 18 & Calibrarea pragului de decizie pe baza scorurilor \\
\texttt{metrics}              & 18 & Calculul și exportul metricilor de evaluare \\
\texttt{mann\_whitney}        & 18 & Testul statistic Mann--Whitney U și mărimea efectului \\
\texttt{plotter\_3d}          & 15 & Vizualizarea tridimensională a clusterelor \\
\texttt{clustering}           & 10 & Clusteringul aglomerativ ierarhic \\
\hline
\textbf{Total}                & \textbf{353} & \\
\hline
\end{tabular}
\end{table}

\subsection{Rezultate}

În urma execuției complete a suitei, toate cele \textbf{353 de teste sunt
trecute cu succes} (\emph{353 passed}), fără eșecuri sau erori. Timpul total de
execuție este de aproximativ 40 de secunde, ceea ce face suita suficient de
rapidă pentru a fi rulată frecvent în timpul dezvoltării. Avertismentele rămase
(18 \emph{warnings}) provin din biblioteci externe (\texttt{numpy},
\texttt{scikit-learn}) și corespund unor situații numerice degenerate
\emph{intenționat} provocate de teste --- de exemplu calculul ariei de sub
curba ROC (AUC) atunci când există o singură clasă în datele de referință ---,
fără a indica defecte ale aplicației.

\paragraph{Defecte identificate și corectate.} Procesul de testare și-a
demonstrat utilitatea prin identificarea unui set de defecte, dintre care cele
mai relevante au fost:

\begin{itemize}
  \item În modulul de testare statistică Mann--Whitney, câmpurile booleene de
        decizie (\texttt{reject\_h0\_alpha05}) erau returnate ca tip boolean
        \texttt{numpy} în loc de tipul boolean nativ Python, ceea ce afecta
        comparațiile stricte de identitate; defectul a fost remediat prin
        conversie explicită de tip.
  \item În testele de extragere a trăsăturilor, indicii care identificau
        trăsăturile \emph{eficiența traiectoriei} și \emph{rata de pauză} nu mai
        corespundeau structurii actualizate a vectorului de trăsături (care a
        fost extins cu descriptori suplimentari precum tremorul, jerk-ul și
        asimetria direcțională); indicii au fost actualizați la pozițiile
        corecte (37, respectiv 43).
  \item Mai multe defecte erau cauzate de \emph{interferența} dintre module de
        test care modificau global tabela de module a interpretorului
        (\texttt{sys.modules}), substituind biblioteci reale cu obiecte
        simulate; izolarea acestor substituiri a eliminat dependențele de ordine
        dintre teste și a restabilit determinismul suitei.
\end{itemize}

Aceste corecții ilustrează rolul dublu al testării: pe de o parte validarea
codului de producție, pe de altă parte menținerea calității și a izolării
propriei infrastructuri de testare.

\subsection{Reproductibilitate și limitări}

Suita de teste este complet deterministă: generatoarele de numere aleatoare sunt
inițializate cu semințe fixe (\texttt{seed}), astfel încât rezultatele sunt
identice la fiecare rulare și independente de mediul de execuție. Lansarea se
face printr-o singură comandă (\texttt{pytest tests/}), iar dependențele sunt
cele declarate în mediul virtual al proiectului.

Ca limitări, suita se concentrează pe testarea unitară și de proprietate, nu
include teste de integrare \emph{end-to-end} care să parcurgă întregul flux de
prelucrare pe date reale, și nu cuantifică explicit acoperirea de cod
(\emph{code coverage}) printr-un instrument dedicat. Aceste direcții constituie
posibile extinderi viitoare, alături de adăugarea unor teste de performanță și a
integrării într-un sistem de integrare continuă (CI) care să ruleze automat
suita la fiecare modificare a codului.

% =====================================================================

\section{Testarea automată a sistemului}
\label{sec:testare}

\subsection{Rolul și obiectivele testării}

Testarea automată reprezintă o componentă esențială a procesului de dezvoltare
a aplicației, având rolul de a verifica în mod sistematic și reproductibil
corectitudinea funcțională a fiecărei unități de cod. Într-un sistem de
prelucrare a datelor și de analiză comportamentală, în care rezultatele finale
(clasificarea participanților, estimarea gradului de dizabilitate, metricile de
evaluare) depind de un lanț lung de transformări numerice, o singură eroare
introdusă într-o funcție de bază se poate propaga silențios până la concluziile
științifice. Suita de teste are tocmai rolul de a interpune o barieră de
verificare automată între implementare și utilizare.

Obiectivele urmărite prin construirea suitei de teste au fost următoarele:

\begin{itemize}
  \item \textbf{Verificarea corectitudinii funcționale} a fiecărei funcții
        publice din modulele aplicației, pe baza unor date de intrare cu
        rezultat cunoscut a priori.
  \item \textbf{Protejarea împotriva regresiilor}: orice modificare ulterioară a
        codului poate fi validată prin re-rularea suitei, asigurând că
        funcționalitățile deja corecte nu sunt deteriorate.
  \item \textbf{Documentarea executabilă a comportamentului așteptat}: testele
        descriu, prin exemple concrete, contractul fiecărei funcții (intrări
        valide, valori de frontieră, cazuri degenerate), constituind o formă de
        specificație vie a codului.
  \item \textbf{Validarea robusteții} în fața cazurilor limită (date lipsă,
        traiectorii cu un singur cadru, distribuții identice, clase unice),
        care apar frecvent în datele experimentale reale.
  \item \textbf{Izolarea defectelor}: în cazul unei erori, granularitatea fină a
        testelor permite localizarea rapidă a modulului și a funcției
        responsabile.
\end{itemize}

\subsection{Metodologia de testare}

Strategia adoptată este aceea a \emph{testării unitare} (unit testing),
completată de elemente de \emph{testare comportamentală} și de verificare a
cazurilor de frontieră. Fiecare test urmează tiparul \emph{Arrange--Act--Assert}
(pregătirea datelor de intrare, apelul funcției testate, verificarea
rezultatului prin aserțiuni), ceea ce asigură lizibilitate și o corespondență
clară între intenția testului și implementarea sa.

Categoriile de teste utilizate sunt:

\begin{itemize}
  \item \textbf{Teste de corectitudine numerică} --- verifică egalitatea
        (eventual aproximativă, cu toleranță controlată) între valoarea
        calculată și cea așteptată teoretic; de exemplu, faptul că distanța
        totală parcursă este nulă pentru o traiectorie staționară.
  \item \textbf{Teste de proprietate / invarianți} --- verifică faptul că
        rezultatul respectă constrângeri structurale independente de datele
        concrete: eficiența traiectoriei este întotdeauna în intervalul
        $[0,1]$, procentele aferente grupurilor de clustering însumează
        $100\%$, vectorul de trăsături nu conține valori \texttt{NaN} sau
        infinite.
  \item \textbf{Teste de cazuri limită} (edge cases) --- traiectorii cu un
        singur cadru sau cu număr minim de cadre, distribuții complet
        separate sau identice în testele statistice, seturi de date vide.
  \item \textbf{Teste de regresie} --- introduse explicit în urma identificării
        și corectării unor defecte, pentru a garanta că acestea nu reapar.
\end{itemize}

\subsection{Implementarea suitei de teste}

\paragraph{Cadrul de testare.} Suita este implementată folosind
\texttt{pytest}, cadrul standard de facto pentru testarea în Python, ales pentru
sintaxa concisă a aserțiunilor (utilizarea directă a instrucțiunii
\texttt{assert}), pentru mecanismul de \emph{fixtures} și pentru capacitățile
avansate de colectare și raportare. Pentru comparațiile numerice se utilizează
\texttt{pytest.approx} și funcțiile \texttt{numpy.testing} (de exemplu
\texttt{assert\_array\_almost\_equal}), care permit specificarea unei toleranțe
controlate, necesară din cauza erorilor de rotunjire inerente aritmeticii în
virgulă mobilă.

\paragraph{Organizarea.} Testele sunt grupate în directorul dedicat
\texttt{tests/}, câte un fișier pentru fiecare modul al aplicației, iar în
interiorul fiecărui fișier testele sunt structurate pe clase care corespund
funcțiilor publice testate (de exemplu \texttt{TestComputeMannWhitney},
\texttt{TestPerformAgglomerativeClustering}). Această corespondență directă
între structura testelor și structura codului facilitează navigarea și
mentenanța.

\paragraph{Configurarea mediului.} Fișierul \texttt{conftest.py} configurează
calea de import (\texttt{sys.path}) astfel încât modulele din \texttt{src/} să
poată fi importate de teste, asigurând independența suitei față de modul de
lansare. Generarea graficelor cu \texttt{matplotlib} este comutată pe backend-ul
neinteractiv \texttt{Agg}, pentru ca testele să poată rula într-un mediu fără
interfață grafică (de exemplu pe un server de integrare continuă).

\paragraph{Izolarea dependențelor.} Pentru a testa fiecare unitate independent
de dependențele sale externe sau costisitoare, s-a folosit tehnica
\emph{mocking}-ului (prin \texttt{unittest.mock}). Astfel, funcțiile de
încărcare a datelor de pe disc, modulul de configurare și anumite biblioteci de
vizualizare sunt înlocuite cu obiecte simulate (\emph{mock objects}), care
returnează valori controlate. Aceasta permite testarea logicii unei funcții
fără a depinde de existența fișierelor reale de date VR, de starea sistemului de
fișiere sau de randarea efectivă a graficelor. Operațiunile de scriere a
fișierelor sunt redirecționate către directoare temporare furnizate de
fixture-ul \texttt{tmp\_path}, astfel încât testele să nu producă efecte
secundare persistente.

\paragraph{Acoperirea.} Suita cuprinde \textbf{353 de teste} care acoperă toate
modulele funcționale ale aplicației. Tabelul~\ref{tab:teste} prezintă
distribuția testelor pe module, împreună cu responsabilitatea fiecărui modul
testat.

\begin{table}[ht]
\centering
\caption{Distribuția testelor pe module}
\label{tab:teste}
\small
\begin{tabular}{l c p{6.5cm}}
\hline
\textbf{Modul testat} & \textbf{Nr. teste} & \textbf{Responsabilitate} \\
\hline
\texttt{optimize\_config}     & 55 & Optimizarea pragurilor și a ponderilor de decizie, încărcarea scorurilor, actualizarea configurației \\
\texttt{file\_management}     & 49 & Gruparea fișierelor pe participant, agregarea trăsăturilor, gestionarea fișierelor de ieșire \\
\texttt{disability\_engine}   & 40 & Determinarea numărului optim de clustere, estimarea statusului de dizabilitate, praguri derivate din date, analiza per participant \\
\texttt{feature\_extractor}   & 26 & Extragerea vectorului de trăsături comportamentale din traiectorie \\
\texttt{scenario\_comparison} & 22 & Compararea și vizualizarea scenariilor pe participanți \\
\texttt{loocv}                & 21 & Validarea încrucișată leave-one-out, praguri Youden și F1 \\
\texttt{ground\_truth\_handler} & 21 & Încărcarea și sincronizarea etichetelor de referință \\
\texttt{reports}              & 20 & Generarea rapoartelor vizuale și a dendrogramelor \\
\texttt{data\_loader}         & 20 & Parsarea vectorilor și încărcarea înregistrărilor VR \\
\texttt{threshold\_calibrator}& 18 & Calibrarea pragului de decizie pe baza scorurilor \\
\texttt{metrics}              & 18 & Calculul și exportul metricilor de evaluare \\
\texttt{mann\_whitney}        & 18 & Testul statistic Mann--Whitney U și mărimea efectului \\
\texttt{plotter\_3d}          & 15 & Vizualizarea tridimensională a clusterelor \\
\texttt{clustering}           & 10 & Clusteringul aglomerativ ierarhic \\
\hline
\textbf{Total}                & \textbf{353} & \\
\hline
\end{tabular}
\end{table}

\subsection{Rezultate}

În urma execuției complete a suitei, toate cele \textbf{353 de teste sunt
trecute cu succes} (\emph{353 passed}), fără eșecuri sau erori. Timpul total de
execuție este de aproximativ 40 de secunde, ceea ce face suita suficient de
rapidă pentru a fi rulată frecvent în timpul dezvoltării. Avertismentele rămase
(18 \emph{warnings}) provin din biblioteci externe (\texttt{numpy},
\texttt{scikit-learn}) și corespund unor situații numerice degenerate
\emph{intenționat} provocate de teste --- de exemplu calculul ariei de sub
curba ROC (AUC) atunci când există o singură clasă în datele de referință ---,
fără a indica defecte ale aplicației.

\paragraph{Defecte identificate și corectate.} Procesul de testare și-a
demonstrat utilitatea prin identificarea unui set de defecte, dintre care cele
mai relevante au fost:

\begin{itemize}
  \item În modulul de testare statistică Mann--Whitney, câmpurile booleene de
        decizie (\texttt{reject\_h0\_alpha05}) erau returnate ca tip boolean
        \texttt{numpy} în loc de tipul boolean nativ Python, ceea ce afecta
        comparațiile stricte de identitate; defectul a fost remediat prin
        conversie explicită de tip.
  \item În testele de extragere a trăsăturilor, indicii care identificau
        trăsăturile \emph{eficiența traiectoriei} și \emph{rata de pauză} nu mai
        corespundeau structurii actualizate a vectorului de trăsături (care a
        fost extins cu descriptori suplimentari precum tremorul, jerk-ul și
        asimetria direcțională); indicii au fost actualizați la pozițiile
        corecte (37, respectiv 43).
  \item Mai multe defecte erau cauzate de \emph{interferența} dintre module de
        test care modificau global tabela de module a interpretorului
        (\texttt{sys.modules}), substituind biblioteci reale cu obiecte
        simulate; izolarea acestor substituiri a eliminat dependențele de ordine
        dintre teste și a restabilit determinismul suitei.
\end{itemize}

Aceste corecții ilustrează rolul dublu al testării: pe de o parte validarea
codului de producție, pe de altă parte menținerea calității și a izolării
propriei infrastructuri de testare.

\subsection{Reproductibilitate și limitări}

Suita de teste este complet deterministă: generatoarele de numere aleatoare sunt
inițializate cu semințe fixe (\texttt{seed}), astfel încât rezultatele sunt
identice la fiecare rulare și independente de mediul de execuție. Lansarea se
face printr-o singură comandă (\texttt{pytest tests/}), iar dependențele sunt
cele declarate în mediul virtual al proiectului.

Ca limitări, suita se concentrează pe testarea unitară și de proprietate, nu
include teste de integrare \emph{end-to-end} care să parcurgă întregul flux de
prelucrare pe date reale, și nu cuantifică explicit acoperirea de cod
(\emph{code coverage}) printr-un instrument dedicat. Aceste direcții constituie
posibile extinderi viitoare, alături de adăugarea unor teste de performanță și a
integrării într-un sistem de integrare continuă (CI) care să ruleze automat
suita la fiecare modificare a codului.
